% ============================================================
\chapter{Compactness in Non-Hausdorff Spaces}
\label{ch:compactness}
% ============================================================

Compactness is one of the most reassuring notions in classical topology: in a Hausdorff space, a compact set is closed, bounded (in a metric), and possesses all the finiteness properties one could wish for. But what happens when one abandons the Hausdorff axiom? Compact sets cease to be closed, limit points are no longer unique, and geometric intuition falters. This is no accident: in domain theory, the spaces that model computations are typically $T_0$ but not $T_1$. The notion of \emph{saturated compact set} --- a compact set equal to the intersection of its open neighbourhoods --- then replaces that of closed compact set, and Smyth compactification plays the role that Alexandroff played in the Hausdorff setting. This chapter develops this refined theory.

\begin{intuition}
In Hausdorff spaces, compactness is well understood: compact sets are closed, intersections
of compact sets are compact, and the Alexandroff one-point compactification yields a Hausdorff
space if and only if the original space is locally compact and Hausdorff. In non-Hausdorff
spaces, compactness presents fundamental subtleties. Compact sets need not be closed, and one
must resort to the notion of \emph{saturated compact set} to recover reasonable behavior.
This chapter develops the theory of compactness adapted to $T_0$ spaces, with applications
to domain theory.
\end{intuition}

% ------------------------------------------------------------
\section{Saturated Compact Sets}
% ------------------------------------------------------------

\begin{definition}[Saturated set]
\index{saturated set}
Let $(X,\tau)$ be a topological space. A subset $A \subseteq X$ is \emph{saturated} if it
equals the intersection of all open sets containing it:
\[
  A = \bigcap \{ U \in \tau : A \subseteq U \}.
\]
Equivalently, $A$ is saturated if and only if $A$ is an upper set for the specialization
order: if $x \in A$ and $x \leq y$ (i.e., $x \in \overline{\{y\}}$), then $y \in A$.
\end{definition}

\begin{remark}
In a $T_1$ space, every subset is saturated since the specialization order is equality.
The notion of saturation is thus only relevant in non-$T_1$ spaces.
\end{remark}

\begin{proposition}[Saturation of a set]
For any subset $A$ of a topological space $X$, the \emph{saturation} of $A$ is
\[
  \uparrow\! A = \{ y \in X : \exists x \in A, \; x \leq y \}
\]
where $\leq$ denotes the specialization order. It is the smallest saturated set
containing $A$.
\end{proposition}

\begin{definition}[Saturated compact set]
\index{saturated compact set}
A subset $Q$ of a topological space is a \emph{saturated compact set} if it is both
compact and saturated.
\end{definition}

\begin{attention}
In a non-Hausdorff space, a compact set need not be closed. For example, in the
Sierpi\'nski space $\mathbb{S} = \{0,1\}$ with $\tau = \{\emptyset, \{1\}, \{0,1\}\}$,
the singleton $\{1\}$ is compact (finite) and open, but not closed. However, it is
saturated.
\end{attention}

\begin{theorem}[Saturated compact sets and intersections]
\label{thm:inter-compact-saturated}
Let $X$ be a topological space. The intersection of a saturated compact set $Q$ and a
closed compact set $K$ is compact. More generally, the finite intersection of saturated
compact sets is a saturated compact set.
\end{theorem}

\begin{proof}
Let $(U_i)_{i \in I}$ be an open cover of $Q_1 \cap Q_2$, where $Q_1$ and $Q_2$ are two
saturated compact sets. For every $x \in Q_1 \setminus Q_2$, since $Q_2$ is saturated,
there exists an open set $V_x$ containing $Q_2$ but not $x$. Then
$(U_i)_{i \in I} \cup (X \setminus Q_2)$ covers $Q_1$, and by compactness of $Q_1$ we
extract a finite subcover. The open sets from this subcover that are among the $U_i$
cover $Q_1 \cap Q_2$.
\end{proof}

\begin{example}
Consider $\N$ equipped with the Alexandrov topology associated with the usual order.
The open sets are the upper sets $\{n, n+1, n+2, \ldots\}$ and $\emptyset$. Every finite
set $\{n\}$ is compact but not saturated (since $\uparrow\!\{n\} = \{n, n+1, \ldots\}$).
\end{example}

% ------------------------------------------------------------
\section{Alexandroff One-Point Compactification for Non-Hausdorff Spaces}
% ------------------------------------------------------------

\begin{definition}[Alexandroff compactification]
\index{Alexandroff compactification}
Let $(X,\tau)$ be a topological space. The \emph{one-point compactification} (or
Alexandroff compactification) is the space $X^* = X \cup \{\infty\}$ equipped with
the topology
\[
  \tau^* = \tau \cup \{ X^* \setminus K : K \subseteq X \text{ compact and saturated} \}.
\]
\end{definition}

\begin{theorem}[Properties of the Alexandroff compactification]
Let $X$ be a topological space.
\begin{enumerate}
  \item $X^*$ is compact.
  \item The inclusion $X \hookrightarrow X^*$ is a topological (open) embedding.
  \item If $X$ is $T_0$, then $X^*$ is $T_0$.
  \item If $X$ is sober, then $X^*$ is sober if and only if $X$ is locally compact.
\end{enumerate}
\end{theorem}

\begin{proof}
(1) Let $(U_i)_{i \in I}$ be an open cover of $X^*$. One of the $U_i$ contains $\infty$,
so it is of the form $X^* \setminus K$ for some saturated compact set $K$. The remaining
open sets cover $K$, and by compactness of $K$ we extract a finite subcover.

(2) The open sets of $X$ in $X^*$ are exactly the elements of $\tau$, since
$U \cap X = U$ for $U \in \tau$.

(3) If $X$ is $T_0$ and $x \in X$, then $\{x\}$ and $\{\infty\}$ are separated by the
open set $X$ of $X^*$.

(4) Omitted (see the next section on locally compact sober spaces).
\end{proof}

\begin{remark}
Unlike the Hausdorff case, the Alexandroff compactification of a non-Hausdorff space is
never Hausdorff (unless the original space is already Hausdorff and locally compact).
Nevertheless, it preserves $T_0$ separation, which is the natural level of separation
in domain theory.
\end{remark}

% ------------------------------------------------------------
\section{Locally Compact Sober Spaces}
% ------------------------------------------------------------

\begin{definition}[Locally compact space]
\index{locally compact space}
A topological space $X$ is \emph{locally compact} if for every $x \in X$ and every open
set $U$ containing $x$, there exist an open set $V$ and a saturated compact set $Q$ such
that $x \in V \subseteq Q \subseteq U$.
\end{definition}

\begin{proposition}
In a sober space, local compactness is equivalent to the following condition: the lattice
of open sets $\mathcal{O}(X)$ is a continuous lattice.
\end{proposition}

\begin{definition}[Way-below relation]
\index{way-below relation}
In a complete lattice $L$, we say that $a$ is \emph{way below} $b$, written $a \ll b$,
if for every directed family $D$ such that $b \leq \bigvee D$, there exists $d \in D$
with $a \leq d$.
\end{definition}

\begin{theorem}[Characterization of locally compact sober spaces]
\label{thm:sober-lc}
For a topological space $X$, the following conditions are equivalent:
\begin{enumerate}
  \item $X$ is sober and locally compact.
  \item The lattice $\mathcal{O}(X)$ is a continuous lattice.
  \item $X$ is homeomorphic to the spectrum of a continuous lattice.
\end{enumerate}
\end{theorem}

\begin{proof}
$(1) \Rightarrow (2)$: If $X$ is locally compact, for every open set $U$ and every
$x \in U$, there exist an open $V$ and a saturated compact $Q$ with
$x \in V \subseteq Q \subseteq U$. One verifies that $V \ll U$ in $\mathcal{O}(X)$
and therefore $U = \bigvee \{V : V \ll U\}$.

$(2) \Rightarrow (3)$: By the representation theorem for continuous lattices, the
spectrum of $\mathcal{O}(X)$ is sober and locally compact, and $\mathcal{O}(X)$ is
isomorphic to its lattice of open sets.

$(3) \Rightarrow (1)$: The spectrum of a continuous lattice is sober by definition and
locally compact because the way-below relation provides the required saturated compact
neighborhoods.
\end{proof}

% ------------------------------------------------------------
\section{Stably Compact Spaces}
% ------------------------------------------------------------

\begin{definition}[Stably compact space]
\index{stably compact space}
A topological space $X$ is \emph{stably compact} if it satisfies:
\begin{enumerate}
  \item $X$ is $T_0$;
  \item $X$ is compact;
  \item $X$ is locally compact;
  \item $X$ is \emph{coherent}: the intersection of any two saturated compact sets is a
    saturated compact set;
  \item $X$ is sober.
\end{enumerate}
\end{definition}

\begin{intuition}
Stably compact spaces are the non-Hausdorff analogues of compact Hausdorff spaces. They
arise naturally as spaces of models in denotational semantics (Scott domains) and in logic
(generalized Stone spaces).
\end{intuition}

\begin{example}[Spectrum of a distributive lattice]
If $L$ is a bounded distributive lattice, its spectrum (the set of prime ideals equipped
with the Hull-Kernel topology) is a stably compact space.
\end{example}

\begin{theorem}[Characterization via frames]
A $T_0$ space is stably compact if and only if its lattice of open sets $\mathcal{O}(X)$
is a continuous lattice and $X$ is compact and coherent.
\end{theorem}

\begin{proposition}[Stability properties]
The class of stably compact spaces is closed under:
\begin{enumerate}
  \item finite products;
  \item closed saturated subspaces;
  \item passage to the patch topology.
\end{enumerate}
\end{proposition}

% ------------------------------------------------------------
\section{The Patch Topology}
% ------------------------------------------------------------

\begin{definition}[Patch (or Lawson) topology]
\index{patch topology}
\index{Lawson topology}
Let $X$ be a topological space. The \emph{patch topology} on $X$ is the topology generated
by the open sets of $X$ and the complements of saturated compact sets of $X$. We write
$X^{\mathrm{patch}}$ for the resulting space.
\end{definition}

\begin{theorem}[Patch of a stably compact space]
\label{thm:patch}
If $X$ is stably compact, then:
\begin{enumerate}
  \item $X^{\mathrm{patch}}$ is compact Hausdorff.
  \item The open sets of $X$ are exactly the open sets of $X^{\mathrm{patch}}$ that are
    upper sets for the specialization order of $X$.
  \item The specialization order of $X$ is a closed partial order in
    $(X^{\mathrm{patch}})^2$.
\end{enumerate}
\end{theorem}

\begin{proof}
(1) The patch topology is finer than the original topology, hence separated (since
saturated compact sets allow separation of points). Compactness follows from Alexander's
subbasis theorem applied to the subbasis consisting of open sets and complements of
saturated compact sets, using coherence of $X$.

(2) and (3) follow from the construction of the patch topology and the properties of the
specialization order.
\end{proof}

\begin{example}
The interval $[0,1]$ equipped with the Scott topology (open sets: $\emptyset$ and the
intervals $(a, 1]$ for $a \in [0,1)$, plus $[0,1]$) is stably compact. Its patch topology
is the usual topology on $[0,1]$.
\end{example}

% ------------------------------------------------------------
\section{De Groot Duality}
% ------------------------------------------------------------

\begin{definition}[Co-compact (de Groot dual) topology]
\index{de Groot duality}
Let $X$ be a topological space. The \emph{de Groot dual topology} $\tau^d$ is the topology
generated by the complements of saturated compact sets of $(X,\tau)$:
\[
  \tau^d = \langle \{ X \setminus Q : Q \text{ saturated compact in } (X,\tau) \} \rangle.
\]
We write $X^d = (X, \tau^d)$ for the dual space.
\end{definition}

\begin{theorem}[De Groot duality for stably compact spaces]
\label{thm:de-groot}
If $X$ is stably compact, then:
\begin{enumerate}
  \item $X^d$ is stably compact.
  \item $(X^d)^d = X$ (the duality is involutive).
  \item The specialization order of $X^d$ is the reverse of the specialization order of $X$.
  \item The patch topologies of $X$ and $X^d$ coincide.
\end{enumerate}
\end{theorem}

\begin{proof}
(1) The saturated compact sets of $X^d$ are exactly the closed sets of $X$ that are lower
sets for the specialization order. One verifies the five axioms one by one.

(2) Follows from the fact that the patch topology determines both $\tau$ and $\tau^d$.

(3) The specialization order of $\tau^d$ is determined by: $x \leq^d y$ iff for every
saturated compact set $Q$ of $\tau$, $y \in Q \Rightarrow x \in Q$, which corresponds to
$y \leq x$ in the specialization order of $\tau$.

(4) The patch topology of $X^d$ is generated by $\tau^d$ and complements of saturated
compact sets of $X^d$, which are open sets of $X$. Hence it coincides with the patch
topology of $X$.
\end{proof}

\begin{keyformulas}[title=\textbf{Summary: de Groot duality}]
\begin{center}
\renewcommand{\arraystretch}{1.3}
\begin{tabular}{|l|c|c|}
\hline
& \textbf{Space} $X$ & \textbf{Dual} $X^d$ \\
\hline
Open sets & $\tau$ & Generated by $X \setminus Q$, $Q$ saturated compact \\
Spec.\ order & $\leq$ & $\geq$ \\
Patch topology & $\tau^{\mathrm{patch}}$ & $(\tau^d)^{\mathrm{patch}} = \tau^{\mathrm{patch}}$ \\
\hline
\end{tabular}
\end{center}
\end{keyformulas}

% ------------------------------------------------------------
\section{Applications in Domain Theory}
% ------------------------------------------------------------

\begin{theorem}[Domains as stably compact spaces]
Let $D$ be a continuous bounded-complete domain (i.e., a continuous dcpo in which every
bounded subset has a supremum). Then $D$ equipped with the Scott topology is stably compact.
\end{theorem}

\begin{corollary}
The interval domain $\mathbf{IR} = \{[a,b] : a \leq b, \; a,b \in \R\}$ ordered by
reverse inclusion is stably compact for the Scott topology.
\end{corollary}

\begin{exercise}
Show that in a $T_0$ space, a compact set is saturated if and only if it is an intersection
of open sets.
\end{exercise}

\begin{exercise}
Let $X$ be a locally compact sober space. Show that for every open set $U$ of $X$, $U$ is
the union of the interiors of saturated compact sets contained in $U$.
\end{exercise}

\begin{exercise}
Let $P$ be a finite ordered set equipped with the Alexandrov topology. Determine the
saturated compact sets of $P$. Deduce the patch topology and the de Groot dual topology.
\end{exercise}

\begin{exercise}
Show that if $X$ is stably compact and $f\colon X \to Y$ is continuous, surjective, and
open, then $Y$ is stably compact.
\end{exercise}

\begin{exercise}[Generalized Stone duality]
Let $L$ be a bounded distributive lattice and $X = \mathrm{Spec}(L)$ its spectrum. Show
that $X$ is stably compact and that $L$ is isomorphic to the lattice of open saturated
compact sets of $X$.
\end{exercise}

\begin{exercise}
Let $X$ be a stably compact space and $Q_1, Q_2$ two saturated compact sets. Show that
$Q_1 \cup Q_2$ and $Q_1 \cap Q_2$ are saturated compact sets, and that the set of
saturated compact sets of $X$ forms a distributive lattice.
\end{exercise}
